\section{金融体系、激励与市场价格}\label{ux91d1ux878dux4f53ux7cfbux6fc0ux52b1ux4e0eux5e02ux573aux4ef7ux683c}

\subsection{The Financial System（金融体系）}\label{the-financial-systemux91d1ux878dux4f53ux7cfb}

金融体系由 markets、financial intermediaries、financial service firms 及其他制度组成，为家庭、企业和政府的金融决策提供环境。它既约束决策者，也扩展其跨期交易和风险分担能力。

金融市场可以有明确交易场所，也可以通过计算机和通信网络形成 \textbf{over-the-counter, OTC} 市场。金融中介以提供金融产品和服务为主营业务，例如银行、投资公司和保险公司。

\subsubsection{Flow of funds}\label{flow-of-funds}

资金从 \textbf{surplus units} 流向 \textbf{deficit units}，两类单位都可能是家庭、企业或政府：

\begin{enumerate}
\def\labelenumi{\arabic{enumi}.}
\tightlist
\item
  资金盈余者与资金短缺者直接缔约；
\item
  通过证券市场完成 direct finance；
\item
  通过银行等中介完成 indirect finance。
\end{enumerate}

\textbf{Disintermediation} 指绕过传统中介、直接或经市场融资，并不表示金融体系中不再需要信息、清算、信用评估与合约执行等功能。

\subsection{Functional Perspective（功能视角）}\label{functional-perspectiveux529fux80fdux89c6ux89d2}

金融机构会随时间和法域改变，而金融功能相对稳定。因此应先问``需要完成什么功能''，再问``由哪种机构完成''。课件列出六项核心功能：

\begin{enumerate}
\def\labelenumi{\arabic{enumi}.}
\tightlist
\item
  跨时间和空间转移资源；
\item
  管理风险；
\item
  清算和结算支付，促进商品、服务与资产交易；
\item
  汇集资源以支持大规模、不可分割的项目，并把权利拆分成可持有的份额；
\item
  提供价格信息，协调分散决策；
\item
  处理信息不对称和激励问题。
\end{enumerate}

课堂材料将货币制度、国际货币体系、金融科技与科技金融纳入功能视角：制度形式会演化，但跨期支付、价值储藏、融资、风险转移和信任生产等底层需求持续存在。

\subsection{Incentive Problems（激励问题）}\label{incentive-problemsux6fc0ux52b1ux95eeux9898}

\begin{longtable}[]{@{}
  >{\raggedright\arraybackslash}p{(\linewidth - 6\tabcolsep) * \real{0.2500}}
  >{\raggedright\arraybackslash}p{(\linewidth - 6\tabcolsep) * \real{0.2500}}
  >{\raggedright\arraybackslash}p{(\linewidth - 6\tabcolsep) * \real{0.2500}}
  >{\raggedright\arraybackslash}p{(\linewidth - 6\tabcolsep) * \real{0.2500}}@{}}
\toprule\noalign{}
\begin{minipage}[b]{\linewidth}\raggedright
问题
\end{minipage} & \begin{minipage}[b]{\linewidth}\raggedright
发生阶段
\end{minipage} & \begin{minipage}[b]{\linewidth}\raggedright
核心机制
\end{minipage} & \begin{minipage}[b]{\linewidth}\raggedright
例子
\end{minipage} \\
\midrule\noalign{}
\endhead
\bottomrule\noalign{}
\endlastfoot
Adverse selection & 签约前 & 高风险者更积极进入交易，交易对手难以识别类型 & 高风险投保人更愿意购买保险 \\
Moral hazard & 签约后 & 受到保障或取得资金后，行为变得更冒险或更少预防损失 & 投保后降低防损努力 \\
Principal--agent problem & 委托关系中 & 代理人目标或信息与委托人不一致 & 经理人追求私人利益而非股东价值 \\
\end{longtable}

一个运作良好的金融体系不会消除这些问题，而是通过筛选、抵押、契约限制、监督、声誉、资本要求和激励相容机制降低其成本。

\subsection{Financial Innovation and Derivatives（金融创新与衍生品）}\label{financial-innovation-and-derivativesux91d1ux878dux521bux65b0ux4e0eux884dux751fux54c1}

金融创新往往由分散的企业家与企业为满足需求、规避约束或利用新技术而产生，而不是完全由中央机构设计。评价创新不能只看技术新颖性，还要看它是否改善某项金融功能、如何重新分配风险，以及是否产生新的激励问题。

\textbf{Forward contract} 约定一方在未来特定日期以预先确定的价格买入资产，另一方必须卖出。\textbf{Futures contract} 是在有组织交易所交易的标准化远期合约。两者都能锁定未来交易价格，但在标准化、交易场所、保证金与结算安排上不同。

\subsection{Financial Market Rates（金融市场价格）}\label{financial-market-ratesux91d1ux878dux5e02ux573aux4ef7ux683c}

本节关注三类价格信号：interest rates、exchange rates 和 stock-market indicators。

\subsubsection{Interest rate}\label{interest-rate}

利率是以某一计价单位表示的承诺收益率，受以下因素影响：

\begin{itemize}
\tightlist
\item
  \textbf{Unit of account}：现金流以哪种货币或商品篮子计价；
\item
  \textbf{Maturity}：距离偿还还有多长时间；
\item
  \textbf{Default risk}：借款人不能或不愿偿还的可能性。
\end{itemize}

一国政府债券在本币意义下可能接近无违约风险，但对外国投资者并非无风险，因为其本币回报还取决于汇率。\textbf{Risk-free} 也不等于零利率。

\subsection{Exchange-Rate Risk Example（汇率风险例题）}\label{exchange-rate-risk-exampleux6c47ux7387ux98ceux9669ux4f8bux9898}

设中国投资者拥有 \(15{,}000\) 元，投资期限为一年：

\[
S_0=8.3\ \text{CNY/USD},\qquad
r_{\text{CNY}}=3\%,\qquad
r_{\text{USD}}=4\%.
\]

直接投资人民币债券，一年后得到

\[
15{,}000(1.03)=15{,}450\ \text{CNY}.
\]

若先换成美元并投资美国债券，则初始美元本金和期末美元价值分别为

\[
\frac{15{,}000}{8.3}=1{,}807.229\ \text{USD},
\]

\[
1{,}807.229(1.04)=1{,}879.518\ \text{USD}.
\]

设一年后的即期汇率为 \(S_1\)（CNY/USD），换回人民币后的收益率为

\[
R_{\text{CNY from USD}}
=\frac{S_1}{S_0}(1+r_{\text{USD}})-1.
\]

\subsubsection{\texorpdfstring{Scenario A：人民币升值，\(S_1=7.2\)}{Scenario A：人民币升值，S\_1=7.2}}\label{scenario-aux4ebaux6c11ux5e01ux5347ux503cs_17.2}

\[
V_1=1{,}879.518\times7.2\approx13{,}532.53\ \text{CNY},
\]

\[
R\approx\frac{13{,}532.53}{15{,}000}-1=-9.78\%.
\]

美元债券本身收益率为正，但人民币升值造成的换汇损失超过了利息收入。

\begin{quote}
{[}!warning{]} 课件数值核对
CH2 p.27 图中写作 \(13{,}533.53\) 元；按该页给定的 \(15{,}000/8.3\times1.04\times7.2\) 计算应约为 \(13{,}532.53\) 元。收益率 \(-9.78\%\) 不受这一小误差影响。
\end{quote}

\subsubsection{\texorpdfstring{Scenario B：人民币贬值，\(S_1=8.7\)}{Scenario B：人民币贬值，S\_1=8.7}}\label{scenario-bux4ebaux6c11ux5e01ux8d2cux503cs_18.7}

\[
V_1=1{,}879.518\times8.7\approx16{,}351.81\ \text{CNY},
\]

\[
R\approx\frac{16{,}351.81}{15{,}000}-1=9.01\%.
\]

可见，未套期保值的外国债券本币回报由外国利率和汇率变动共同决定。

\subsubsection{Forward exchange rate and no-arbitrage}\label{forward-exchange-rate-and-no-arbitrage}

若用一年期远期合约锁定换回人民币的汇率 \(F_{0,1}\)，要使人民币投资与``换美元---买美元债券---远期换回人民币''具有相同期末价值，应满足

\[
15{,}000(1+r_{\text{CNY}})
=\frac{15{,}000}{S_0}(1+r_{\text{USD}})F_{0,1}.
\]

因此

\[
F_{0,1}
=S_0\frac{1+r_{\text{CNY}}}{1+r_{\text{USD}}}
=8.3\times\frac{1.03}{1.04}
\approx8.22\ \text{CNY/USD}.
\]

这就是 \textbf{covered interest parity} 的一期间形式。远期合约把未知的 \(S_1\) 替换为签约时确定的 \(F_{0,1}\)，从而锁定本币回报；若市场远期汇率显著偏离该值且不存在交易成本、信用风险与融资限制，就会出现套利压力。

\subsection{Today's Takeaways}\label{todays-takeaways}

\begin{enumerate}
\def\labelenumi{\arabic{enumi}.}
\tightlist
\item
  金融学的统一问题是：如何跨时间配置资源，并在不确定性下评价现金流。
\item
  家庭与企业面对相似的消费、投资、融资和风险管理问题，但企业还必须处理资本结构、控制权和代理冲突。
\item
  机构形式会变化，金融体系的六项基本功能更稳定；功能视角有助于理解金融创新。
\item
  金融契约不仅转移资金，也重新分配风险、信息权与控制权。
\item
  外国无违约风险债券对本国投资者仍有汇率风险；外国利率高于本国利率并不保证本币收益更高。
\item
  远期汇率的无套利定价来自两个确定性投资路径期末价值相等，而不是对未来即期汇率的主观预测。
\end{enumerate}

\subsection{Terms}\label{terms}

\begin{longtable}[]{@{}
  >{\raggedright\arraybackslash}p{(\linewidth - 4\tabcolsep) * \real{0.3333}}
  >{\raggedright\arraybackslash}p{(\linewidth - 4\tabcolsep) * \real{0.3333}}
  >{\raggedright\arraybackslash}p{(\linewidth - 4\tabcolsep) * \real{0.3333}}@{}}
\toprule\noalign{}
\begin{minipage}[b]{\linewidth}\raggedright
Term
\end{minipage} & \begin{minipage}[b]{\linewidth}\raggedright
中文
\end{minipage} & \begin{minipage}[b]{\linewidth}\raggedright
简要含义
\end{minipage} \\
\midrule\noalign{}
\endhead
\bottomrule\noalign{}
\endlastfoot
Financial system & 金融体系 & 支持金融契约、资产交易与风险转移的市场和机构集合 \\
Asset allocation & 资产配置 & 在不同资产类别之间分配储蓄 \\
Capital budgeting & 资本预算 & 识别、评价并实施长期投资项目 \\
Capital structure & 资本结构 & 企业价值在股权、债务等资本来源之间的配置 \\
Working capital & 营运资本 & 流动资产与流动负债之差 \\
Agency conflict & 代理冲突 & 委托人与代理人的目标或信息不一致 \\
Disintermediation & 脱媒 & 融资绕过传统金融中介 \\
Adverse selection & 逆向选择 & 签约前由隐藏类型导致的选择偏差 \\
Moral hazard & 道德风险 & 签约后由隐藏行为导致的风险上升 \\
Forward contract & 远期合约 & 约定未来按固定价格交易资产的非标准化合约 \\
Unit of account & 计价单位 & 合约支付所采用的货币或价值尺度 \\
Covered interest parity & 抛补利率平价 & 经远期套保后的跨币种无套利利率关系 \\
\end{longtable}
